\chapter{Sistemas Procedurais}

A premissa fundamental do Void Descent é que \emph{nenhum asset é importado}. Texturas,
geometria, áudio, comportamento de inimigos e mapas são gerados algoritmicamente a cada
\emph{run}. Este capítulo descreve os três sistemas procedurais centrais usando
\textbf{Sequence Diagrams}, que representam interação ordenada entre componentes ao
longo do tempo, o tipo UML adequado para descrever \emph{processos} multi-etapas.

\section{Geração de \emph{dungeon}}

\diagfig{diagrams/dungeon-gen.png}{Diagrama de Sequência da geração procedural de um andar}{dungeon-gen}

\subsection{Algoritmo de \emph{room placement}}

A função \texttt{dungeon\_generate(floor\_id)} segue cinco etapas:

\begin{enumerate}
  \item \textbf{Seeding}. \texttt{random\_seed(0xBADC0DE1 + floor\_id)} reinicia o LCG.
    Andares com mesmo \emph{seed} produzem o mesmo mapa, garantindo determinismo do
    \emph{layout} dentro de uma \emph{run}. Isso é importante para que
    \texttt{count\_enemies\_in\_room} identifique consistentemente a mesma sala em
    \emph{frames} diferentes.

  \item \textbf{Posicionamento de salas}. Em um \emph{loop} com até 500 tentativas,
    são sorteados retângulos com $w, h \in [7, 11]$ tiles e $x, y$ válidos no mapa
    $64 \times 64$. Cada candidato é rejeitado se sobrepõe (com \emph{gap} mínimo de
    1 tile) qualquer sala já aceita. O número alvo é $4 + \mathit{floor\_id}$ salas,
    exceto no Andar 3, que tem 5.

  \item \textbf{Tipagem}. \texttt{rooms[0]} é a sala inicial; a última é
    \texttt{chest} (Andar 1 e 2) ou \texttt{boss} (Andar 3). No Andar 3 a segunda
    sala também é \texttt{chest}, garantindo equipamento antes do confronto final.

  \item \textbf{Carving} dos corredores. Para cada par de salas adjacentes na lista,
    um corredor em forma de L é escavado conectando os centroides, com escolha
    aleatória da ordem (horizontal-primeiro ou vertical-primeiro).

  \item \textbf{Construção da geometria}. Cada tile é convertido em chamada a
    \texttt{push\_floor} (chão) ou \texttt{push\_block} (parede). O resultado é
    \emph{baked} no \emph{vertex buffer} do WebGL. Não há atualização incremental
    durante o \emph{gameplay}, o que permite que \texttt{level\_num\_verts} marque
    o fim da geometria estática.
\end{enumerate}

\subsection{Validação}

O posicionamento por descarte tem chance teórica de gerar mapas onde uma sala fique
ilhada. Na prática, com 500 tentativas e o algoritmo de conexão sequencial
\texttt{i $\to$ i+1}, todas as salas ficam alcançáveis a partir da inicial via grafo
linear de corredores.

\section{Síntese de áudio}

\diagfig{diagrams/audio-synth.png}{Diagrama de Sequência do \emph{pipeline} de síntese de áudio}{audio-synth}

\subsection{Pipeline}

A \texttt{audio\_init} é assíncrona:

\begin{enumerate}
  \item Cria o \texttt{AudioContext}, inicialmente \emph{suspended} pela política de
    \emph{autoplay}, e resumido na primeira interação do usuário.
  \item Cria \texttt{audio\_master\_gain} (um \texttt{GainNode}) e o conecta ao
    \texttt{destination}. Todas as fontes posteriores passam por este nó,
    permitindo \emph{mute} unificado via
    \texttt{audio\_master\_gain.gain.value = 0}.
  \item Dispara a síntese assíncrona da música. \texttt{sonantxr\_generate\_song}
    processa cada \emph{track} em \emph{chunks} de 33\,ms, intercalando com
    \texttt{setTimeout(0)} para não bloquear a \emph{main thread} excessivamente.
  \item Em paralelo, dispara a síntese de oito SFX (\emph{shoot}, \emph{hit},
    \emph{hurt}, \emph{beep}, \emph{pickup}, \emph{terminal}, \emph{explode},
    \emph{dash}). Cada um é síncrono, mas curto.
\end{enumerate}

\subsection{Estrutura interna do sintetizador}

Cada nota processada em \texttt{SoundGenerator.\_genSound} executa as seguintes
etapas, sobre os $\sim$10\,000 \emph{samples} de duração da nota:

\begin{itemize}
  \item Dois osciladores (\texttt{osc1}, \texttt{osc2}) com formas selecionáveis
    (\textit{sin}, \textit{square}, \textit{saw}, \textit{tri}) e \emph{detune}
    aplicável.
  \item Envelope ADSR (\texttt{attack}, \texttt{sustain}, \texttt{release})
    calculado por \emph{sample}.
  \item Filtro biquad de estado configurável (\textit{lopass}, \textit{hipass},
    \textit{bandpass}, \textit{notch}) com ressonância.
  \item LFO modulando a frequência do filtro.
  \item \emph{Pan} e \emph{delay} aplicados em pós-processamento sobre o
    \texttt{Float32Array} resultante.
\end{itemize}

A faixa ambiente atual usa um único \emph{track} de \emph{drone} com
$rowLen = 11025$ e $songLen = 16$\,s, gerando um \emph{loop} sem cortes audíveis.

\section{Atlas de tiles e \emph{tinting}}

O atlas de tiles é embarcado no código como \emph{string} \texttt{base64}
(\texttt{TILES\_B64} em \texttt{game.js}). No \emph{boot} de cada andar,
\texttt{tint\_atlas} aplica um \emph{hue shift} ao atlas inteiro pixel-a-pixel,
produzindo a paleta característica de cada andar: ciano para o Andar 1
($\Delta H = 0{,}55$), laranja para o Andar 2 ($\Delta H = 0{,}15$) e magenta
para o Andar 3 ($\Delta H = 0{,}75$). O resultado é uma textura
$\mathit{tile\_atlas\_size}^2$ vinculada ao GL via \texttt{renderer\_bind\_image}.

A conversão de espaço é \texttt{RGB} $\rightarrow$ \texttt{HSL} $\rightarrow$
\texttt{shifted HSL} $\rightarrow$ \texttt{RGB}, com saturação amplificada em 40\%
para preservar o \emph{look} \emph{neon} sobre fundo preto.
